\chapter{Complexity — ND Combinators}
%%% generated by the blueprint pipeline from TCSlib.Complexity.TuringMachine.Build.NDCombinators
\section{Overview}

Declarations of the Lean module \texttt{TCSlib.Complexity.TuringMachine.Build.NDCombinators}.

\section{Declarations}

\begin{theorem}[Bounds on a sign]
\lean{Turing.signType_bounds}
\label{Turing.signType_bounds}
\leanok
\difficulty{1}
Every sign $s\in\{-1,0,1\}$ satisfies $-1\le s\le 1$.
\end{theorem}

\begin{theorem}[Step from a specified state]
\lean{Turing.NDTM.stepWith_of_state}
\label{Turing.NDTM.stepWith_of_state}
\leanok
\uses{Turing.Action.apply, Turing.Cfg, Turing.Cfg.inputSymbol, Turing.Cfg.workTapeSymbols, Turing.NDTM, Turing.NDTM.stepWith}
\difficulty{3}
Let $M$ be a nondeterministic machine with $k$ work tapes, let $b$ be a choice bit, and
let $c$ be a configuration on a fixed input. If the state component of $c$ is $q$, then
one step under choice $b$ is obtained by applying to $c$ the action prescribed by the
transition function at $q$, using the input symbol and work-tape symbols scanned in $c$.
\end{theorem}

\begin{definition}[Agreement under tape shifts]
\lean{Turing.NDTM.AgreeShift}
\label{Turing.NDTM.AgreeShift}
\leanok
\uses{Turing.Cfg}
Let $c$ and $c'$ be configurations of the same $k$-tape machine on the same input, let
$r$ be a natural number, and let $d$ assign an integer shift to each tape. They agree
under the shifts $d$ beyond distance $r$ when they have the same state, input-head
position, and output; the head on tape $i$ in $c'$ is shifted by $d(i)$ from its
position in $c$; and for every tape $i$ and every cell $z$ at or to the right of the
head in $c$ minus $r$, the contents of cell $z+d(i)$ in $c'$ equal those of cell $z$ in
$c$.
\end{definition}

\begin{theorem}[Shifted agreement after one choice]
\lean{Turing.NDTM.AgreeShift.stepWith}
\label{Turing.NDTM.AgreeShift.stepWith}
\leanok
\uses{Turing.Action.apply, Turing.Cfg, Turing.Cfg.inputSymbol, Turing.Cfg.workTapeSymbols, Turing.NDTM, Turing.NDTM.AgreeShift, Turing.NDTM.stepWith, Turing.NDTM.stepWith_of_halt, Turing.signType_bounds}
\difficulty{6}
Let $M$ be a nondeterministic multi-tape Turing machine, let $b$ be a choice bit, let
$r$ be a natural number, let $d$ assign an integer shift to each work-tape index, and
let $c,c'$ be configurations that agree with shift $d$ and radius $r+1$. Then the
configurations reached by one step under $b$ from $c$ and $c'$ agree with shift $d$ and
radius $r$.
\end{theorem}

\begin{theorem}[Shifted agreement under a choice word]
\lean{Turing.NDTM.AgreeShift.runWith}
\label{Turing.NDTM.AgreeShift.runWith}
\leanok
\uses{Turing.Cfg, Turing.NDTM, Turing.NDTM.AgreeShift, Turing.NDTM.AgreeShift.stepWith, Turing.NDTM.runWith, Turing.NDTM.runWith_cons}
\difficulty{4}
Let $M$ be a nondeterministic multi-tape Turing machine, let $w$ be a finite binary
choice word, let $r$ be a natural number, let $d$ assign an integer shift to each
work-tape index, and let $c,c'$ be configurations that agree with shift $d$ and radius
$r+|w|$. Then the configurations reached by running $M$ from $c$ and $c'$ under $w$
agree with shift $d$ and radius $r$.
\end{theorem}

\begin{theorem}[Reflexivity of shifted agreement]
\lean{Turing.NDTM.AgreeShift.refl}
\label{Turing.NDTM.AgreeShift.refl}
\leanok
\uses{Turing.Cfg, Turing.NDTM.AgreeShift}
\difficulty{1}
Let $r$ be a natural number and $c$ a configuration. The configurations $c$ and $c$
agree with shift $0$ and radius $r$: their states, input-head positions, and outputs
coincide; their work-tape head positions differ by $0$ on every tape; and for every tape
$i$ and cell $z$ with the head position on tape $i$ minus $r$ at most $z$, the tape
contents at $z$ coincide.
\end{theorem}

\begin{theorem}[Work-tape displacement during a run]
\lean{Turing.NDTM.abs_runWith_workTapePos_sub_le}
\label{Turing.NDTM.abs_runWith_workTapePos_sub_le}
\leanok
\uses{Turing.Action.apply, Turing.Cfg, Turing.Cfg.inputSymbol, Turing.Cfg.workTapeSymbols, Turing.NDTM, Turing.NDTM.runWith, Turing.NDTM.runWith_cons, Turing.NDTM.stepWith, Turing.NDTM.stepWith_of_halt, Turing.NDTM.stepWith_of_state, Turing.signType_bounds}
\difficulty{5}
Let $M$ be a nondeterministic multi-tape Turing machine, $w$ a finite binary choice
word, $c$ a configuration, and $i$ a work-tape index. The absolute difference between
the position of head $i$ after running $M$ from $c$ under $w$ and its initial position
is at most $|w|$.
\end{theorem}

\begin{theorem}[Unvisited work-tape cells are unchanged]
\lean{Turing.NDTM.runWith_workTapes_of_le}
\label{Turing.NDTM.runWith_workTapes_of_le}
\leanok
\uses{Turing.Action.apply, Turing.Cfg, Turing.Cfg.inputSymbol, Turing.Cfg.workTapeSymbols, Turing.NDTM, Turing.NDTM.abs_runWith_workTapePos_sub_le, Turing.NDTM.runWith, Turing.NDTM.runWith_cons, Turing.NDTM.runWith_nil, Turing.NDTM.stepWith, Turing.NDTM.stepWith_of_halt, Turing.NDTM.stepWith_of_state}
\difficulty{5}
Let $M$ be a nondeterministic multi-tape Turing machine, let $w$ be a choice word, and
let $c$ be a configuration on a fixed input. For every work tape $i$ and every integer
cell position $z$, if the initial head position on tape $i$ plus $|w|$ is at most $z$,
then the content of cell $z$ on tape $i$ after running under $w$ equals its initial
content.
\end{theorem}

\begin{definition}[Mapping a machine's output]
\lean{Turing.NDTM.mapOutput}
\label{Turing.NDTM.mapOutput}
\leanok
\uses{Turing.NDTM}
Let $M$ be a nondeterministic $k$-tape machine and let $f$ map symbols to optional
symbols. The machine obtained by mapping its output through $f$ keeps the same initial
state and, in each transition, keeps every part of the action unchanged except its
optional output symbol, which is replaced by the result of applying $f$ to it and
flattening the optional value.
\end{definition}

\begin{definition}[Mapping configuration output]
\lean{Turing.NDTM.mapOutputCfg}
\label{Turing.NDTM.mapOutputCfg}
\leanok
\uses{Turing.Cfg}
Let $c$ be a configuration of a $k$-tape machine, and let $f$ map symbols to optional
symbols. The mapped configuration keeps the state, input-head position, work-tape
contents, and work-head positions of $c$; its output is obtained by applying $f$ to each
output symbol in order and omitting the symbols for which $f$ returns no value.
\end{definition}

\begin{theorem}[Output Mapping Commutes with One Step]
\lean{Turing.NDTM.mapOutput_stepWith}
\label{Turing.NDTM.mapOutput_stepWith}
\leanok
\uses{Turing.Action.apply, Turing.Cfg, Turing.Cfg.inputSymbol, Turing.Cfg.workTapeSymbols, Turing.NDTM, Turing.NDTM.mapOutput, Turing.NDTM.mapOutputCfg, Turing.NDTM.stepWith, Turing.NDTM.stepWith_of_halt}
\difficulty{4}
Let $M$ be a nondeterministic Turing machine, let $f$ be a partial symbol map, let $b$
be a choice bit, and let $c$ be a configuration. Taking one step under $b$ in the
machine whose transition outputs are mapped by $f$, from the configuration whose output
is mapped by $f$, yields the configuration obtained by taking one step in $M$ from $c$
under $b$ and then mapping its output by $f$.
\end{theorem}

\begin{theorem}[Output Mapping Commutes with a Run]
\lean{Turing.NDTM.mapOutput_runWith}
\label{Turing.NDTM.mapOutput_runWith}
\leanok
\uses{Turing.Cfg, Turing.NDTM, Turing.NDTM.mapOutput, Turing.NDTM.mapOutputCfg, Turing.NDTM.mapOutput_stepWith, Turing.NDTM.runWith, Turing.NDTM.runWith_cons}
\difficulty{3}
Let $M$ be a nondeterministic Turing machine, let $f$ be a partial symbol map, let $w$
be a finite word of choice bits, and let $c$ be a configuration of $M$. Running the
machine with its transition outputs mapped by $f$, starting from the configuration whose
output is mapped by $f$, gives the same result as first running $M$ from $c$ under $w$
and then mapping the resulting output by $f$.
\end{theorem}

\begin{theorem}[Output Mapping Preserves the Initial Configuration]
\lean{Turing.NDTM.mapOutput_initCfg}
\label{Turing.NDTM.mapOutput_initCfg}
\leanok
\uses{Turing.NDTM, Turing.NDTM.initCfg, Turing.NDTM.mapOutput, Turing.NDTM.mapOutputCfg}
\difficulty{1}
Let $M$ be a nondeterministic Turing machine with initial state $q_0$, and let $f$ be a
partial symbol map. For every input word $x$, mapping the output symbols of the initial
configuration of $M$ on $x$ by $f$ gives the same configuration as the initial
configuration of the machine whose transition outputs are mapped by $f$.
\end{theorem}

\begin{definition}[Mapping a finite machine's output]
\lean{Turing.FinNDTM.mapOutput}
\label{Turing.FinNDTM.mapOutput}
\leanok
\uses{Turing.FinNDTM, Turing.NDTM.mapOutput}
Let $N$ be a finite nondeterministic machine over the Boolean alphabet, and let $f$ map
Boolean symbols to optional Boolean symbols. The mapped machine has the same number of
work tapes and the same state set as $N$, and uses the transition rule obtained by
mapping each transition's optional output through $f$ while leaving its other action
data unchanged.
\end{definition}

\begin{definition}[Left-branch action embedding]
\lean{Turing.FinNDTM.coinBranchAct₀}
\leanok
\uses{Turing.Action}
Let $N_0$ and $N_1$ be finite nondeterministic machines over the Boolean alphabet, and
let $a$ be an action of $N_0$. The corresponding action on the combined $k_0+k_1$ work
tapes keeps $a$'s input movement and output, uses $a$'s work-tape actions on the first
$k_0$ tapes and leaves the remaining $k_1$ tapes unwritten and stationary. A successor
state of $a$ is sent to the left component of the disjoint union of the two state sets,
with the result wrapped as a nonhalting state; if $a$ has no successor state, the
combined action also halts.
\end{definition}

\begin{definition}[Action in the second machine branch]
\lean{Turing.FinNDTM.coinBranchAct₁}
\leanok
\uses{Turing.Action}
Let $N_0$ and $N_1$ be finite nondeterministic Turing machines over the Boolean
alphabet, and let $a$ be a one-step action of $N_1$. The corresponding action on the
combined machine has the same input-head movement and output; it leaves every work tape
assigned to $N_0$ unchanged, uses $a$ on the work tapes assigned to $N_1$, and maps any
successor state of $a$ into the second summand of the combined state type, with the
outer option indicating a nonhalting state.
\end{definition}

\begin{definition}[Combined machine with a choice branch]
\lean{Turing.FinNDTM.coinBranch}
\label{Turing.FinNDTM.coinBranch}
\leanok
\uses{Turing.FinNDTM, Turing.FinNDTM.coinBranchAct₀, Turing.FinNDTM.coinBranchAct₁}
Let $N_0$ and $N_1$ be finite nondeterministic Turing machines over the Boolean
alphabet. Their combined machine has the sum of their work-tape counts and state type
$	ext{Option}(N_0.	ext{State} \oplus N_1.	ext{State})$, with initial state none. From
that state, a choice bit selects the initial state of $N_0$ or $N_1$; thereafter, in
either branch, the machine applies the selected machine's transition for that bit to its
own work tapes, leaving the other machine's tapes unchanged and embedding successor
states in the corresponding summand. The transition out of the initial state leaves the
input and work-tape heads stationary, makes no writes or output, and enters the selected
initial state.
\end{definition}

\begin{definition}[Configuration embedded in the first branch]
\lean{Turing.FinNDTM.coinBranchCfg₀}
\leanok
\uses{Turing.Cfg, Turing.FinNDTM.coinBranch}
Let $N_0$ and $N_1$ be finite nondeterministic Turing machines over the Boolean
alphabet, and let $x$ be a Boolean input string. A configuration of $N_0$ on $x$ is sent
to a configuration of their combined machine on the same input by embedding its state in
the first summand (preserving a halting state), retaining its input-head position and
output, and placing its work tapes and head positions in the first block. Every tape in
the second block is blank at every position and has its head at position $0$.
\end{definition}

\begin{definition}[Configuration embedded in the second branch]
\lean{Turing.FinNDTM.coinBranchCfg₁}
\leanok
\uses{Turing.Cfg, Turing.FinNDTM.coinBranch}
Let $N_0$ and $N_1$ be finite nondeterministic Turing machines over the Boolean
alphabet, and let $x$ be a Boolean input string. A configuration of $N_1$ on $x$ is sent
to a configuration of their combined machine on the same input by embedding its state in
the second summand (preserving a halting state), retaining its input-head position and
output, and placing its work tapes and head positions in the second block. Every tape in
the first block is blank at every position and has its head at position $0$.
\end{definition}

\begin{theorem}[Step preserves the first branch]
\lean{Turing.FinNDTM.coinBranch_stepWith₀}
\leanok
\uses{Turing.Action.apply, Turing.Cfg, Turing.Cfg.inputSymbol, Turing.Cfg.workTapeSymbols, Turing.FinNDTM.coinBranch, Turing.FinNDTM.coinBranchAct₀, Turing.FinNDTM.coinBranchCfg₀, Turing.NDTM.stepWith, Turing.NDTM.stepWith_of_halt, Turing.NDTM.stepWith_of_state}
\difficulty{5}
Let $N_0,N_1$ be finite nondeterministic Turing machines with Boolean symbols, let $x$
be a Boolean input, let $b\in\{0,1\}$, and let $c$ be a configuration of $N_0$ on $x$.
Taking one step of the combined machine under choice $b$ from the configuration embedded
in its first branch gives the embedded configuration obtained by taking that step of
$N_0$ from $c$.
\end{theorem}

\begin{theorem}[Step preserves the second branch]
\lean{Turing.FinNDTM.coinBranch_stepWith₁}
\leanok
\uses{Turing.Action.apply, Turing.Cfg, Turing.Cfg.inputSymbol, Turing.Cfg.workTapeSymbols, Turing.FinNDTM.coinBranch, Turing.FinNDTM.coinBranchAct₁, Turing.FinNDTM.coinBranchCfg₁, Turing.NDTM.stepWith, Turing.NDTM.stepWith_of_halt, Turing.NDTM.stepWith_of_state}
\difficulty{5}
Let $N_0,N_1$ be finite nondeterministic Turing machines with Boolean symbols, let $x$
be a Boolean input, let $b\in\{0,1\}$, and let $c$ be a configuration of $N_1$ on $x$.
Taking one step of the combined machine under choice $b$ from the configuration embedded
in its second branch gives the embedded configuration obtained by taking that step of
$N_1$ from $c$.
\end{theorem}

\begin{theorem}[Run preserves the first branch]
\lean{Turing.FinNDTM.coinBranch_runWith₀}
\leanok
\uses{Turing.Cfg, Turing.FinNDTM.coinBranch, Turing.FinNDTM.coinBranchCfg₀, Turing.FinNDTM.coinBranch_stepWith₀, Turing.NDTM.runWith, Turing.NDTM.runWith_cons}
\difficulty{4}
Let $N_0,N_1$ be finite nondeterministic Turing machines with Boolean symbols, let $x$
be a Boolean input, let $w$ be a choice word, and let $c$ be a configuration of $N_0$ on
$x$. Running the combined machine that simulates $N_0$ in its first branch under $w$,
starting from the configuration embedded in that branch, yields the embedded
configuration reached by running $N_0$ from $c$ under $w$.
\end{theorem}

\begin{theorem}[Run preserves the second branch]
\lean{Turing.FinNDTM.coinBranch_runWith₁}
\leanok
\uses{Turing.Cfg, Turing.FinNDTM.coinBranch, Turing.FinNDTM.coinBranchCfg₁, Turing.FinNDTM.coinBranch_stepWith₁, Turing.NDTM.runWith, Turing.NDTM.runWith_cons}
\difficulty{4}
Let $N_0,N_1$ be finite nondeterministic Turing machines with Boolean symbols, let $x$
be a Boolean input, let $w$ be a choice word, and let $c$ be a configuration of $N_1$ on
$x$. Running the combined machine that simulates $N_1$ in its second branch under $w$,
starting from the configuration embedded in that branch, yields the embedded
configuration reached by running $N_1$ from $c$ under $w$.
\end{theorem}

\begin{theorem}[Initial step selects the coin branch]
\lean{Turing.FinNDTM.coinBranch_stepWith_init}
\label{Turing.FinNDTM.coinBranch_stepWith_init}
\leanok
\uses{Turing.Cfg.init, Turing.FinNDTM.coinBranch, Turing.FinNDTM.coinBranchCfg₀, Turing.FinNDTM.coinBranchCfg₁, Turing.NDTM.initCfg, Turing.NDTM.stepWith}
\difficulty{5}
Let $N_0,N_1$ be finite nondeterministic Turing machines with Boolean symbols, let $x$
be a Boolean input, and let $b\in\{0,1\}$. Taking one step of the combined machine from
its initial configuration on $x$ using choice $b$ gives the embedded initial
configuration of $N_1$ when $b=1$, and the embedded initial configuration of $N_0$ when
$b=0$.
\end{theorem}

\begin{theorem}[Run of the selected coin branch]
\lean{Turing.FinNDTM.coinBranch_runWith}
\label{Turing.FinNDTM.coinBranch_runWith}
\leanok
\uses{Turing.FinNDTM.coinBranch, Turing.FinNDTM.coinBranchCfg₀, Turing.FinNDTM.coinBranchCfg₁, Turing.FinNDTM.coinBranch_runWith₀, Turing.FinNDTM.coinBranch_runWith₁, Turing.FinNDTM.coinBranch_stepWith_init, Turing.NDTM.initCfg, Turing.NDTM.runWith, Turing.NDTM.runWith_cons}
\difficulty{4}
Let $N_0,N_1$ be finite nondeterministic Turing machines with Boolean symbols, let $x$
be a Boolean input, let $b\in\{0,1\}$, and let $w$ be a choice word. Starting from its
initial configuration on $x$, the machine that first selects between $N_0$ and $N_1$
according to $b$ and then simulates the selected machine under $w$ reaches the
configuration obtained by embedding that machine's run into the corresponding branch: it
is the embedded run of $N_1$ when $b=1$ and the embedded run of $N_0$ when $b=0$.
\end{theorem}
